\section{Integrals Reducible to Rational Functions}

It is important to know that the antiderivative of any rational function can be
expressed using the methods from the previous section. Similarly, it is useful
to know that there are other cases that can be reduced to the integration of a
rational function.

\emph{Rational functions in $e^x$.}
If the integrand function $f(x)$ can be expressed in the form
\[
  f(x) = R(e^{\lambda x})
\]
where $R$ is a rational function and $\lambda\neq 0$,
then the integral can be solved through the substitution $t = e^{\lambda x}$.
Indeed, we have
\[
\begin{cases}
 e^{\lambda x} = t\\
 x = \frac{\ln t}{\lambda}\\
 dx = \frac{1}{\lambda t}\, dt
 \end{cases}
\]
and the integrand function becomes a rational function:
\[
 \int R(e^{\lambda x})\, dx = \Enclose{\int \frac{R(t)}{\lambda t}\, dt}_{t=e^{\lambda x}}
\]

\begin{example}
Let us solve the integral
\[
  \int \frac{2\sqrt{e^x} + e^{2x}}{e^x-4}\, dx.
\]
\end{example}
\begin{proof}[Solution.]
We express the integrand function in terms of $e^{\frac x 2}$:
\[
  \frac{2\sqrt{e^x} + e^{2x}}{e^x-4}
  =\frac{2e^{\frac x 2}+e^{4\frac x 2}}{e^{2\frac x 2}-4}.
\]
By making the substitution $t=e^{\frac x 2}$, $x=2\ln t$, $dx=\frac 2 t\, dt$
we obtain a rational function in $t$:
\[
  \int \frac{2\sqrt{e^x} + e^{2x}}{e^x-4}\, dx
  = \int \frac{2t + t^4}{t^2-4}\cdot \frac 2 t dt
  = 2\int \frac{2+t^3}{t^2-4}\, dt.
\]
By performing polynomial division and partial fraction decomposition, we obtain
\[
 \int 2t\, dt + \int \frac{5}{t-2}\, dt + \int \frac{3}{t+2}\, dt
 = t^2 + 5\ln\abs{t-2} + 3\ln\abs{t+2}
\]
and thus, substituting $t=e^{\frac x 2}$, we obtain the result
\[
 e^x + 5 \ln \abs{\sqrt{e^x}-2} + 3 \ln\enclose{\sqrt{e^x}+2}.
\]
\end{proof}

\emph{Rational functions in $\sin^2 x$, $\cos^2 x$, $\sin x \cdot \cos x$.}
If the integrand function $f(x)$ can be expressed in the form
\[
  f(x) = R(\sin^2 x, \cos^2 x, \sin x \cdot \cos x)
\]
where $R$ is a rational function (i.e., a ratio of polynomials in the three
indicated variables),
then the integral can be solved through the substitution $t=\tg x$. Indeed,
observing that
\[
  1 + \tg^2 x = 1 + \frac{\sin^2 x}{\cos^2 x} = \frac{1}{\cos^2 x}
\]
from which
\begin{align*}
\cos^2 x &= \frac{1}{1+\tg^2 x}\\
\sin^2 x &= \tg^2 x \cdot\cos^2 x\\
\sin x \cdot \cos x &= \tg x \cdot \cos^2 x
\end{align*}
by setting $t = \tg x$ we have:
\begin{equation}\label{eq:466324}
\begin{cases}
  \cos^2 x = \frac{1}{1+t^2}\\
  \sin^2 x = \frac{t^2}{1+t^2}\\
  \sin x \cdot \cos x = \frac{t}{1+t^2}\\
  dx = \frac{1}{1+t^2}\, dt
\end{cases}
\end{equation}
and the integrand function becomes a rational function.

\begin{example}
\label{ex:35663}
Let us calculate
\[
  \int \frac{1}{\cos x \cdot ( \sin x + \cos x)}\, dx.
\]
\end{example}
\begin{proof}[Solution.]
The integrand function can be expressed in the form
\[
  \frac{1}{\sin x \cdot \cos x + \cos^2x}.
\]
By performing the substitution~\eqref{eq:466324}
we obtain
\[
  \int \frac{1}{\frac t {1+t^2} + \frac{1}{1+t^2}}\cdot \frac{1}{1+t^2}\, dt
  = \int \frac{1}{t+1}\, dt = \ln \abs{t+1} = \ln \abs{\tg x + 1}.
\]
\end{proof}

\emph{More generally, rational functions of $\sin x$ and $\cos x$.}
If the integrand function $f(x)$ can be expressed in the form
\[
  f(x) = R(\sin x, \cos x)
\]
where $R$ is a rational function, then the integral can be solved
through the substitution $t=\tg \frac{x}{2}$. Indeed, with this substitution we
have
(using the half-angle formulas and reducing to the previous case)
\begin{equation}\label{eq:3675323}
  \begin{cases}
    \cos x = \cos^2 \frac x 2 - \sin^2 \frac x 2 = \frac{1-t^2}{1+t^2} \\
    \sin x = 2 \sin \frac x 2 \cos \frac x 2 = \frac{2t}{1+t^2}\\
    dx = \frac{2}{1+t^2}\, dt.
  \end{cases}
\end{equation}
Again, with this substitution, the integrand function becomes rational.

\begin{remark}
It should be noted that the substitution~\eqref{eq:3675323} could always be
used in place of~\eqref{eq:466324} as it is more general.
However, usually, if it is possible to use the substitution~\eqref{eq:466324},
the integral then becomes simpler to calculate. Try it with the integral
from example~\ref{ex:35663}!
\end{remark}

\begin{example}
Let us calculate
\[
 \int \frac{1}{\sin x}\, dx.
\]
\end{example}
\begin{proof}[Solution.]
Using the substitution~\eqref{eq:3675323}
we obtain
\begin{align*}
  \int \frac{1}{\sin x}\, dx
  &= \int \frac{1}{\frac{2t}{1+t^2}}\frac{2}{1+t^2}\, dt \\
  &= \int \frac{1}{t}\, dt = \ln\abs t = \ln \abs{\tg \frac x 2}.
\end{align*}
\end{proof}

\emph{Rational functions with radicals.}
If the function $f(x)$ can be expressed in the form:
\[
  f(x) = R(\sqrt[n] x)
\]
where $R$ is a rational function, then the integral can be solved through
the substitution $x = t^n$. Indeed, with this substitution we have
\begin{equation}\label{eq:4675821}
\begin{cases}
  \sqrt[n] x = t\\
  dx = n t^{n-1}\, dt.
\end{cases}
\end{equation}

\begin{example}
Let us calculate
\[
  \int \frac{\sqrt[4]{x}}{\sqrt[2]{x} + \sqrt[3]{x}}\, dx.
\]
\end{example}
\begin{proof}[Solution.]
Observe that the integrand function can be written
as a rational function of $t=\sqrt[12]{x}$ (we have chosen
the least common multiple of the indices involved: $12 =
\mathit{lcm}\ENCLOSE{4,2,3}$)
\[
  \frac{\sqrt[4]{x}}{\sqrt{x} + \sqrt[3]{x}}
  = \frac{\sqrt[12]{x^3}}{\sqrt[12]{x^6} + \sqrt[12]{x^4}}.
\]
Thus, using the substitution \eqref{eq:4675821} with $n=12$, we have
\begin{align*}
\int \frac{t^3}{t^6 + t^4}\cdot 12 t^{11}\, dt
&= 12 \int \frac{t^{14}}{t^6+t^4}\, dt
 = 12 \int \frac{t^{10}}{t^2+1}\, dt.
\end{align*}
Proceeding with polynomial division, we obtain
\begin{align*}
\MoveEqLeft{12 \int \Enclose{t^8 - t^6 + t^4 - t^2 + 1 - \frac{1}{t^2+1}}\, dt} \\
&= 12 \Enclose{\frac{t^9}{9} - \frac{t^7}{7} + \frac{t^5}{t} - \frac{t^3}{3} + t
- \arctg t} \\
&= \frac 4 3 \sqrt[4]{x^3} - \frac{12}{7}\sqrt[12]{x^7}
+ \frac{12}{5}\sqrt[12]{x^5} - 4 \sqrt[4]{x} + 12 \sqrt[12]{x} - 12 \arctg
\sqrt[12]{x}.
\end{align*}
\end{proof}